\chapter{方法}
\label{chap:method}

本章详述自进化强化学习系统的核心方法：多智能体采样架构（\ref{sec:sampling}）、差分驱动的可信奖励（\ref{sec:reward}）、以及策略优化目标（\ref{sec:objective}）。

\section{多智能体采样架构}
\label{sec:sampling}

采样的目标是让策略在真实环境中\textbf{自主产生有信息量的多轮交互经验}。本文用三个协同的智能体完成这一过程：执行者、观察者、提问者。

\subsection{执行者与沙箱环境}

\textbf{执行者}（actor）即被训练的策略 $\pi_\theta$。它在一个\textbf{代码沙箱}中以 ReAct 方式求解任务：每一步生成一段思考并可选地发起一次工具调用（如执行 Python、读写文件、运行命令），沙箱执行后把结果作为观察返回，执行者据此继续，直至给出最终答复或达到步数上限。沙箱提供了一个\textbf{可复现、可取证}的真实环境——任务是否完成，最终要看沙箱里的状态，而不是执行者怎么说。

一次多步 ReAct 交互构成\textbf{一条轨迹}。轨迹中，执行者自己生成的 token 参与策略梯度，工具返回的观察 token 则被掩码，不计入损失。

\subsection{观察者与状态取证}

\textbf{观察者}（observer）的职责是回答"这条轨迹\textbf{实际}改变了环境的什么"。它不依赖执行者的自述，而是在轨迹执行\textbf{前后}对沙箱状态各拍一次快照，计算二者的\textbf{差分}：

\begin{itemize}
  \item \textbf{文件系统差分}：哪些文件被新建、修改、删除，以及变更的\textbf{内容}（对结构化文件如表格、文档，进一步抽取其文本内容进入差分）。
  \item \textbf{系统状态差分}：本轮安装的软件包、开启的端口、启动的进程等。
\end{itemize}

差分是\textbf{确定性}的——它由快照比对得到，不经过大模型的主观归纳。观察者据此产出一份\textbf{状态报告}，作为后续奖励打分的客观依据。这份报告与执行者的\textbf{自我声明}并置：当执行者声称"已把结果写入 12345"而差分显示文件内容实为 99999 时，二者的矛盾被直接暴露。这正是抑制奖励作弊的取证基础。

\subsection{提问者与多轮交互}

真实用户很少一句话说清需求，而是在看到初步结果后\textbf{追问}。为使采样贴近部署分布，\textbf{提问者}（questioner）扮演一个\textbf{具人设}的模拟用户：给定一个种子任务与一种人设（语气、偏好、专业背景各异），它在执行者每轮答复后生成一句符合人设的追问，把单个种子任务扩展为\textbf{一段多轮对话}。当提问者判断需求已被满足时以显式信号结束会话。

这样，一个种子查询不再只产生一条单轮轨迹，而是产生一段包含若干轮追问的多轮交互经验，显著提升了单位采样的信息量，也让训练分布更接近真实的多轮使用场景。采样规模在"查询数 $\times$ 组大小 $\times$ 追问轮数"三个维度上权衡：在固定的总轨迹预算下，减少种子查询数、以固定轮数的多轮追问补足，可在不增加总量的前提下把采样从单轮转为多轮。

\section{差分驱动的可信奖励}
\label{sec:reward}

本文\textbf{不用规则奖励}，而由一个\textbf{外部冻结的大模型裁判}（frozen LLM judge）按统一细则对每条轨迹打分。与常见做法的关键区别在于：裁判打分\textbf{锚定于观察者提供的环境真实状态差分}，执行者的自述仅作交叉核对——从取证机制上抑制奖励作弊。裁判在四个维度打分：

\begin{itemize}
  \item \textbf{Done}（完成度，$0/1$）：任务是否真正完成（以状态差分为准）。
  \item \textbf{Correctness}（正确性，$[0,1]$）：产出是否正确（有验收点时对照，否则由裁判依据差分语义判断）。
  \item \textbf{Trajectory}（轨迹质量，$[0,1]$）：工具调用质量、有无冗余/重复步、推理是否连贯；\textbf{并惩罚"加了不该加的东西"}（装多余的包、改任务范围外的文件等越界行为）。
  \item \textbf{Safety}（安全，\textbf{二元 $0/1$}）：是否出现破坏性/越权操作。作为\textbf{二元门控}而非连续打分——只判"有没有危险操作"这一可判定问题，一致性更高。
\end{itemize}

最终奖励按下式\textbf{乘性}聚合：

\begin{equation}
  r =
  \begin{cases}
    (0.4\cdot \text{correctness} + 0.4\cdot \text{trajectory} + 0.2)\cdot \text{safety}, & \text{if done}=1,\\[4pt]
    0.4\cdot \text{trajectory}\cdot \text{safety}, & \text{if done}=0.
  \end{cases}
  \label{eq:reward}
\end{equation}

安全作为 $\{0,1\}$ 乘性门控：一旦出现危险操作（safety$=0$），无论其余维度多高，整条奖励清零；未完成任务仍给轨迹分（过程有价值），但无完成基线分与正确性分。裁判返回严格 JSON，解析失败重试一次，再失败则丢弃该轨迹；同一 GRPO 组内轨迹丢弃过半则整组丢弃，避免用残缺的组估计优势。

\section{策略优化目标}
\label{sec:objective}

策略更新采用标准 GRPO\upcite{B8}。对每个查询采样一组 $G$ 条轨迹，以组内奖励的均值与标准差归一化得到各轨迹的优势，再以裁剪的策略梯度更新，省去价值网络。本文不引入面向持续学习的额外正则，仅保留一项\textbf{固定开启的熵正则}以预防多轮交互中的策略坍缩（Echo Trap\upcite{B4}）：

\begin{equation}
  L = L_{\text{GRPO}} - \lambda_{\text{ent}}\, \mathcal{H}[\pi_\theta], \qquad \lambda_{\text{ent}}=0.001.
\end{equation}

熵正则在所有实验中固定开启、不参与消融——它不是本文的方法创新，而是保证多轮 GRPO 训练稳定的必要一项。本文训练的目的在于\textbf{验证整套自进化系统能端到端稳定运转}，因此采用标准 GRPO，而非在某个特定算法细节上求最优。
